\documentclass[10pt,a4paper]{amsart}
\usepackage[margin=19mm]{geometry}
\usepackage{amsmath,amssymb,mathtools}
\usepackage[scr=rsfs,cal=euler]{mathalfa}
\usepackage{xfrac}
\usepackage[unicode,hidelinks,bookmarks=false]{hyperref}
\newcommand{\ie}{i.e.\ }
\newcommand{\cf}{cf.\ }
\newcommand{\resp}{respectively\ }
\newcommand{\Subsection}{\subsection}
\providecommand{\nobreakdash}{\nobreak\hbox{-}}
\setlength{\emergencystretch}{2em}
\multlinegap0pt
\usepackage{bm}
\usepackage{bbm}
\usepackage{dsfont}
\usepackage{xspace}
\usepackage{cancel}
\usepackage{mathtools}
\usepackage{amsthm}
\makeatletter
\@ifundefined{theorem}{\newtheorem{theorem}{Theorem}}{}
\@ifundefined{proposition}{\newtheorem{proposition}[theorem]{Proposition}}{}
\makeatother
\DeclareMathOperator{\supp}{supp}
\newcommand{\FIN}{\mathrm{FIN}}
\newcommand{\TS}{\mathrm{TS}}
\title[Ramsey-Theoretic Finiteness in Choiceless Set Theory: The Go]{Ramsey-Theoretic Finiteness in Choiceless Set Theory: The Gower's Collapse and Rainbow and Canonical Separations}
\author{Brinda Venkataramani}
\date{Source: arXiv:2607.22196v1 (CC BY 4.0)}
\begin{document}
\maketitle

\noindent\emph{Source and licence.} Ramsey-Theoretic Finiteness in Choiceless Set Theory: The Gower's Collapse and Rainbow and Canonical Separations. Version arXiv:2607.22196v1 (CC BY 4.0), distributed under the Creative Commons Attribution 4.0 International licence, as recorded in the arXiv licence metadata for this exact version and shown on the abstract page https://arxiv.org/abs/2607.22196v1. Full rights notice: licenses/arxiv-2607.22196-CC-BY-4.0.txt.\par\smallskip

\noindent\textit{Editorial scope.} Counted unit: the Proposition whose source label is \texttt{None} in the article (source0.tex lines 328--330, source label \texttt{None}), delivered together with the article's own complete original proof of it (source0.tex lines 332--372, one \texttt{proof} environment of 41 lines). No mathematics is written, repaired or re-derived: statement and proof are the article's own text line for line. Required context retained uncounted: none. Every definition, display and label of those lines is retained unchanged. Omitted from this package and not redistributed: 1--327, 331--331, 373--837. Nothing inside a retained excerpt is omitted or altered: every retained body is delivered exactly as the article's lines are written, so no omission receipt of any kind accompanies this package. Citations: the retained excerpts carry no citation command at all, so no bibliography entry of the article is transcribed and no reference is redistributed. Reference adaptation: none was needed and none was made; every reference inside the delivered unit resolves inside it, so no unresolved reference or citation is produced. Printed slips kept literal: no printed word, symbol, hypothesis or number of the counted statement or of its proof was altered. Layout and class: the article's own class is \texttt{amsart} with packages including geometry, amsmath, amssymb, amsthm, amsfonts, mathrsfs, enumitem, tikz-cd, hyperref; the wrapper is the lane's pinned \texttt{amsart} editorial wrapper at 10pt on A4, which restates the theorem-like environments the retained text needs and adds the article's own macro definitions that the retained text needs (\texttt{\textbackslash FIN}, \texttt{\textbackslash TS}, \texttt{\textbackslash supp}), with extra wrapper packages (bm, bbm, dsfont, xspace, cancel, mathtools). The delivered numbering is the unit's own. Assets: no retained span contains a figure environment, a graphics-inclusion command, a PDF-page inclusion, a table or a TikZ picture; no image file is copied into this package, the package declares no asset dependency and no image file is redistributed with it. Licence: version arXiv:2607.22196v1 of this article is distributed under the Creative Commons Attribution 4.0 International licence, as recorded in the arXiv licence metadata for this exact version and shown on the abstract page https://arxiv.org/abs/2607.22196v1. A full rights notice is delivered at licenses/arxiv-2607.22196-CC-BY-4.0.txt. Characters: every retained excerpt is pure ASCII, so the delivered engine, XeLaTeX with its default Latin Modern text font, reports no missing character.
\begin{proposition}
Let \(k\geq 1\).  If \([X]^{<\omega}\) is Dedekind-infinite, then \(X\) is \(G_k\)-infinite.
\end{proposition}

\begin{proof}
By the preceding lemma, fix an infinite sequence \((P_j)_{j<\omega}\) of pairwise disjoint nonempty finite subsets of \(X\).

We define a map
\[
    \Phi:\FIN_k(\omega)\to \FIN_k(X)
\]
by spreading the value of a function on \(\omega\) across the corresponding finite block \(P_j\).  Namely, for \(f\in\FIN_k(\omega)\), define \(\Phi(f):X\to\{0,1,\dots,k\}\) by
\[
    \Phi(f)(x)=
    \begin{cases}
        f(j), & x\in P_j,\\
        0, & x\notin \bigcup_{j<\omega}P_j.
    \end{cases}
\]
This is well-defined because the \(P_j\)'s are pairwise disjoint.  Moreover, \(\Phi(f)\in\FIN_k(X)\): the support of \(\Phi(f)\) is contained in the finite union \(\bigcup_{j\in \supp(f)}P_j\), and since \(f\) attains the value \(k\) at some \(j\), the function \(\Phi(f)\) attains the value \(k\) on the nonempty set \(P_j\).

The map \(\Phi\) preserves the tetris operation and finite sums along disjoint supports.  More precisely,
\[
    \Phi(T^r f)=T^r\Phi(f)
\]
for each relevant \(r\), and if \(f_1,\dots,f_s\in\FIN_k(\omega)\) have pairwise disjoint supports, then
\[
    \Phi\left(\sum_{i=1}^s f_i\right)=\sum_{i=1}^s \Phi(f_i).
\]
Since each \(P_j\) is nonempty, \(\Phi\) is injective.  Consequently, \(\Phi\) sends infinite block sequences in \(\FIN_k(\omega)\) to infinite block families in \(\FIN_k(X)\), and it sends tetris sums to the corresponding tetris sums.

Now let
\[
    c:\FIN_k(X)\to 2
\]
be an arbitrary coloring.  Pull it back to a coloring
\[
    d:\FIN_k(\omega)\to 2
\]
by
\[
    d(f)=c(\Phi(f)).
\]
By Gowers' \(\FIN_k\) theorem on \(\omega\), there is an infinite block sequence \((f_i)_{i<\omega}\) in \(\FIN_k(\omega)\) such that \(\TS(\{f_i:i<\omega\})\) is monochromatic for \(d\).  Then \((\Phi(f_i))_{i<\omega}\) is an infinite block family in \(\FIN_k(X)\), and its tetris sums are monochromatic for \(c\).  Since \(c\) was arbitrary, \(X\) is \(G_k\)-infinite.
\end{proof}

\end{document}
